% ctx-prefacio.tex
%
% Copyright (C) 2022--2026 José A. Navarro Ramón <janr.devel@gmail.com>
% Licencia del código GPLv2
% Licencia Creative Commons Recognition Non-Commercial Share-alike.
% (CC-BY-NC-SA)


\startcomponent ctx-prefacio
\product tutorialLuaMetaFun

\startchapter[title={Prefacio}]

Objetivos que pretendo conseguir:
\startitemize[n]
\item Practicar, familiarizarme y clasificar conceptos básicos de \MetaFun en
  \ConTeXt\ \LMTX.
\item Explicar a \emph{mi yo del futuro} los conceptos y ténicas anteriores, para que
  pueda volver a asimilarlas en la medida de lo posible, si lo llegara a necesitar.
\item Practicar y familiarizarme con el sistema tipográfico \ConTeXt.
\stopitemize

Mi experiencia con paquetes gráficos en \TeX\ y derivados se basa casi
exclusivamente en TikZ. Este sistema es muy intuitivo y, sobre todo, tiene un manual
casi perfecto y lo seguiré utilizando, sobre todo en Lua\LaTeX.

Mi interés por \ConTeXt viene de lejos, pero debido a que la documentación está
muy desperdigada y no siempre actualizada, con muchos comandos caducados, me ha
sido muy difícil hacerme con este sistema.

Cuando se integró el lenguaje \type{Lua} en el sistema \ConTeXt, me volví a
interesar y decidí que, aunque no sin dificultad, empezara a utilizarlo documentando
mis avances en forma de código, así estas notas las estoy escribiendo en \ConTeXt\
\LMTX.

Aunque en este sistema también se puede dibujar en TikZ, pretendo utilizar
\MetaPost\ / \MetaFun\ debido a que está más integrado en el sistema, y además,
siempre me ha llamado la atención.

En lugar de documentar \MetaPost\ /\MetaFun en forma de código únicamente como
hago con \ConTeXt, quiero dejar unas notas que me sirvan para asimilar el paquete,
que está integrado en \ConTeXt\ \LMTX. Estas notas se materializan en est texto.

Dicho esto, no pretendo escribir un compendio completo, ni siquiera ordenado,
sobre el sistema \MetaPost\ / \MetaFun. No por nada, sino por falta de tiempo, pues
tengo otras obligaciones e intereses.

Por tanto, la motivación de estas notas ha sido para formar un tutorial que me sirva
a utilizar estas herramientas explicármelo a mí.
Se puede desprender de todo lo anterior que he hecho hincapié en todo aquello que en
algún momento he encontrado más complicado, y no he dedicado ninguna o casi ninguna
atención a detalles que tengo suficientemente asimilados.

Me pregunto si este tutorial podría servir a alguien. No lo sé.
Tampoco pretendo que el tutorial se termine en algún momento. Cada vez que lo repaso
lo modifico algo. El límite es el tiempo del que dispongo, muy limitado, pues tengo
otros intereses que tampoco quiero abandonar, física, matemáticas, música, idiomas,
ajedrez, etc.
En el fondo, como ya estoy jubilado, lo que quiero es disfrutar del viaje, ya que no
sé cuándo se terminará mi billete.

En todo caso, según la licencia, cualquiera que pudiera estar interesado en el texto
y no lo considerara lo suficientemente didáctico o accesible, puede modificarlo y
usarlo para sus propios objetivos. Para ello, consulta la licencia, que he separado
en dos partes:
\startitemize[n]
\item El código fuente, con licencia
  %\href{https://www.gnu.org/licenses/old-licenses/gpl-2.0.html#SEC1}{GPLv2}, se encuentra
  %en \href{https://github.com/janr57/retazosmatematicas}{GitHub}.
\item En cambio, el fichero en formato de texto obtenido, que se puede distribuir por
  separado, tiene licencia
  %\href{https://creativecommons.org/licenses/by-nc-sa/4.0/}{Creative Commons BY-NC-SA}.
\stopitemize

\stopchapter

\stopcomponent


 
%%% Local Variables:
%%% coding: utf-8
%%% mode: context
%%% ConTeXt-Mark-version: "IV"
%%% TeX-engine: luatex
%%% TeX-master: "../ctx-tutorialLuaMetaFun.tex"
%%% End:
